\section{Introduction}
\label{sec:intro}

An aero-engine comes off the wing before it fails, and the decision rests on
an estimate of how many cycles it can still run. That estimate, the remaining
useful life (RUL), is now usually produced by a deep network
\cite{jardine2006review,saxena2008damage,lei2018machinery,hu2022prognostics}.
Physics-based degradation models have largely been replaced: convolutions
pick up local patterns \cite{babu2016deep,li2018remaining}, recurrent layers
track state \cite{zheng2017long,yuan2016fault,xiang2021multicellular}, and
attention encoders handle long dependencies
\cite{wang2024dvgtformer,zhang2024tatfa,kim2025dkam}. Most applied papers now
use all three together, a convolutional front end (often a TCN
\cite{bai2018empirical}), a bidirectional recurrent stage and one or more
self-attention layers \cite{vaswani2017attention}, with differencing,
positional encodings and dropout added along the way. We call this
combination the hybrid backbone.

Accuracy tables on NASA's C-MAPSS benchmark \cite{saxena2008damage} nearly
always come from one training run per configuration. Hyperparameters are
tuned as best-of-$k$, and whether the same configuration would train again
is not reported. We have an archive of 161 runs of one such backbone, and the
omission turns out to matter. On the two subsets with six operating
conditions, 37.2\% of FD002 runs and 24.2\% of FD004 runs finished by
predicting one fixed number for every engine, the mean of the training
labels. None of the 81 runs on single-condition data did. Training loss
falls normally, so a single RMSE near 43.8 looks like a mediocre
hyperparameter choice, but the asymmetric score used in service is
15--250$\times$ worse. Failures also cluster in particular window and depth
settings, and for some of our best configurations the outcome changes with
the random seed. Retraining such a model can quietly replace a good predictor
with one that prints the fleet-average RUL for every engine.

A run that ends this way is easy to recognise once you know what to look
for. Section~\ref{sec:phenomenon} gives five signs that can be read from a
log or a saved checkpoint, one of them a formula for the test RMSE of a
failed run (Eq.~(\ref{eq:fingerprint})) that matches every such run in the
archive to within 0.01 cycles. Why it happens took longer to work out. Probes
on saved checkpoints and per-epoch records from six instrumented runs
(Section~\ref{sec:mechanism}) show that every multi-condition run first
outputs the same constant value: its internal features record which
operating condition the engine is in rather than how far it has worn, and its
attention weights become identical at every time step. Training starts to
improve again only when gradients in the convolutional branch suddenly grow
about a hundredfold, and this has to happen before early stopping ends the
run. Section~\ref{sec:method} turns these observations into a short set of
rules for choosing window length and depth, which we call CAS-Fuse, plus
three cheap checks that can be applied to any finished model.
Section~\ref{sec:experiments} tests the rules over five seeds.

We are not trying to win the accuracy table; specialised models do better on
some subsets (Section~\ref{sec:experiments}). Whether a configuration trains
at all, reliably, when the seed changes has drawn little attention. Among the
hybrid backbones used in applied RUL work, we have not found another
configuration with a multi-seed audit, or an explanation of why its
neighbours sometimes fail.
